\uuid{C19k}
\exo7id{4171}
\titre{exo7 4171}
\auteur{quercia}
\organisation{exo7}
\datecreate{2010-03-11}
\isIndication{false}
\isCorrection{false}
\chapitre{Fonction de plusieurs variables}
\sousChapitre{Dérivée partielle}
\cpge{mpsi-mp2i.fonctions-deux-variables.b, pcsi.fonctions-deux-variables.b, ptsi.fonctions-deux-variables.b, bcpst1.fonctions-deux-variables.c, mp-mpi.calcul-differentiel-optimisation.b, pc.fonctions-plusieurs-variables.c, psi.fonctions-plusieurs-variables.c, pt.fonctions-vectorielles-courbes.b}
\module{Analyse}
\niveau{L2}
\difficulte{}

\contenu{
\question{
\def\arc#1{\mathop{#1}\limits^\frown}
Soit $f : {U \subset \R^n} \to {\R^p}$ de classe $\mathcal{C}^1$ dont les dérivées
partielles sont bornées sur~$U$ et $t\in I\  \mapsto M_t$ une courbe paramétrée dans
$U$ de classe $\mathcal{C}^1$. Pour $a,b\in I$ comparer les longueurs des arcs
$\arc{M_aM_b}$ et $\arc{f(M_a)f(M_b)}$.
}
}
